\section{GC e memória em container}
\label{sec:memoria-container}

\subsection*{A memória que o container vê não é o heap}

O guia \citetitle{doc:quarkus-performance-measure} mede
\termo{RSS}{\emph{Resident Set Size}: memória física de fato ocupada pelo
processo, incluindo heap e memória nativa.}, e não o heap, porque o heap
é só uma parte: a JVM também usa memória nativa para metadados de
classes, pilhas de threads, código compilado e estruturas do GC. É o RSS
que o limite do container corta. Somando o que a Parte~\ref{parte:java-jvm} e o manual do
\texttt{java} (\citetitle{doc:java-cmd}) dizem sobre cada parte:

\small
\begin{longtable}{@{}p{3cm} p{4.4cm} p{4.8cm}@{}}
\toprule
Parte & Opção & Padrão no JDK~25, segundo o manual \\
\midrule \endhead
Heap & \texttt{-XX:MaxRAMPercentage} ou \texttt{-Xmx} & 25\% da memória
  disponível (\texttt{MinRAMPercentage}, 50\%, vale para heaps pequenos,
  de cerca de 125~MB) \\
Metaspace (classes) & \texttt{-XX:MaxMetaspaceSize} & Sem limite \\
Code cache (JIT) & \texttt{-XX:ReservedCodeCacheSize} & 240~MB no
  máximo (48~MB sem compilação em camadas) \\
Pilha por thread & \texttt{-Xss} & 1024~KB em Linux/x64, 2048~KB em
  Linux/AArch64 \\
Buffers diretos & \texttt{-XX:MaxDirectMemorySize} & Escolhido pela JVM
  se não definido \\
\bottomrule
\end{longtable}
\normalsize

O mesmo manual diz que, em Linux, \texttt{UseContainerSupport} vem
ligado: a JVM detecta a memória e as CPUs do container e dimensiona os
recursos com isso. Consequência direta: num container de 1~GB sem
opções, o heap máximo fica em torno de 256~MB, e o resto sobra para
memória nativa, que pode ou não ser suficiente. E as CPUs vistas pela JVM
dimensionam os pools de threads do GC e do \texttt{ForkJoinPool};
\texttt{-XX:ActiveProcessorCount} sobrescreve esse número.

\subsection*{Modo JVM: que GC usar}

A seção~\ref{sec:gc} trouxe a regra do guia de GC da Oracle: deixe a JVM
escolher, ajuste o heap e só troque de coletor se o requisito não for
atendido. Em container, isso vira:
\begin{itemize}
\item Container pequeno, abaixo de duas CPUs ou de 1792~MB: no JDK~25
  a JVM já escolhe o Serial, o perfil que o guia associa a ele
  (\citetitle{doc:gc-ergonomics}). Quem quer o G1 nesse container
  precisa passar \texttt{-XX:+UseG1GC}, e toda medição deve registrar o
  coletor que rodou (\texttt{-Xlog:gc}, seção~\ref{sec:gc}).
\item API com requisito de p99: compare G1 e ZGC no teste de carga,
  com o mesmo limite de memória de produção.
\item Fixe \texttt{-XX:MaxRAMPercentage} (ou \texttt{-Xmx}) explicitamente
  e deixe margem para a memória nativa; ajuste olhando o RSS, não o heap.
\end{itemize}

\subsection*{Modo nativo: outro GC, outro padrão}

O guia \citetitle{doc:graalvm-memory} diz que o Serial GC é o padrão no
Native Image. O G1 existe só no Oracle GraalVM (não no Community
Edition, base do Mandrel), em Linux AMD64 e AArch64. E o detalhe que mais
surpreende: sem heap máximo definido, um executável com Serial GC usa
até 80\% do tamanho da memória física. A página fala em memória física e
não descreve como o limite do container entra nessa conta; por isso,
passe \texttt{-Xmx} em runtime (ou \texttt{-R:MaxHeapSize} no build) e
não dependa do padrão. A mesma página avisa que o RSS pode subir
temporariamente durante uma coleta, o que importa em ambientes com
limite de memória, como containers.

\subsection*{A decisão}

\begin{description}
\item[Quando dar mais memória] RSS próximo do limite sob carga, pausas de
  GC frequentes no JFR, ou heap pequeno demais para o conjunto de dados
  em uso.
\item[Quando reduzir] RSS estável muito abaixo do limite no pico de
  carga, com folga de pausa no p99. Em escala, cada megabyte por réplica
  vira custo.
\item[O que medir] RSS com \texttt{ps -o rss} ou \texttt{pmap -x} (não
  \texttt{docker stats}, segundo o guia de medição) em repouso e no pico
  do teste de carga; pausas e alocação com JFR; reinícios por
  \emph{OOMKilled} no orquestrador. Anote JDK, GC, limites de memória e
  CPU junto de cada número, como pede a seção~\ref{sec:gc}.
\end{description}
